\documentclass[letterpaper,11pt]{article}

\usepackage[empty]{fullpage}
\usepackage[hidelinks]{hyperref}
\usepackage{fontawesome5}
\usepackage{enumitem}
\usepackage{parskip}

% Adjust margins
\addtolength{\oddsidemargin}{-0.4in}
\addtolength{\evensidemargin}{-0.4in}
\addtolength{\textwidth}{0.8in}
\addtolength{\topmargin}{-0.5in}
\addtolength{\textheight}{1.0in}

\ifdefined\pdfgentounicode
  \input{glyphtounicode}
  \pdfgentounicode=1
\fi

\begin{document}

\begin{center}
    {\Huge \scshape Juan Becerra} \\ \vspace{4pt}
    \textbf{\large Software Engineer -- Gameplay, Animation \& Tools} \\ \vspace{4pt}
    \small Seattle, WA ~ $|$ ~
    \raisebox{-0.1\height}\faPhone\ (208) 550-9191 ~ $|$ ~
    \href{mailto:becerrajuan007@gmail.com}{\raisebox{-0.2\height}\faEnvelope\ \underline{becerrajuan007@gmail.com}} ~ $|$ ~
    \href{https://hoowan.dev}{\raisebox{-0.2\height}\faGlobe\ \underline{hoowan.dev}} \\ \vspace{2pt}
    \href{https://linkedin.com/in/hoowan/}{\raisebox{-0.2\height}\faLinkedin\ \underline{linkedin.com/in/hoowan}} ~ $|$ ~
    \href{https://github.com/hoowan-dev}{\raisebox{-0.2\height}\faGithub\ \underline{github.com/hoowan-dev}}
\end{center}

\vspace{4pt}
\hrule
\vspace{14pt}

\noindent October 26, 2026

\vspace{8pt}
\noindent \textbf{Romain Almes} \\
EA Studios - Respawn \\
Electronic Arts \\
Vancouver, BC / Los Angeles, CA

\vspace{8pt}
\noindent \textbf{Re: Gameplay Engineer (Apex Legends) (Role ID: 216048)}

\vspace{8pt}
\noindent Dear Romain Almes,

I am writing to express my strong interest in the \textbf{Gameplay Engineer} position at Respawn 
Entertainment. Over the past five years at Electronic Arts (four years as a Gameplay \& Animation 
Software Engineer at \textbf{Maxis} on \emph{Project Rene} (\emph{The Sims}) and currently at 
\textbf{Lifestyle}), I have focused on building responsive, high-performance player-facing gameplay 
systems in C++ while partnering closely with designers, animators, QA, and production. I would be 
thrilled to bring my passion for fluid movement and moment-to-moment gameplay feel to Apex Legends.

During my tenure at \textbf{EA Maxis Studios}, I implemented core player-facing locomotion and 
runtime animation systems in Unreal Engine 5 and C++, including motion-matching, direct/indirect 
control schemes, and custom IK solvers, to deliver responsive, tactile character controls.
I also architected a multi-threaded, asynchronous path-planning and navigation API that 
enabled deterministic routing and real-time obstacle avoidance for shared multiplayer simulation. 
Working across a large-scale multi-disciplinary codebase, I partnered directly with gameplay 
designers and technical animators to prototype, iterate, and polish gameplay features while 
managing technical risk and long-term maintainability. To support rapid design iteration and 
stability, I authored in-editor tuning tools and visual debuggers and profiled runtime memory 
and compute across character evaluation graphs and spatial queries.

In my current role at \textbf{EA Lifestyle}, I have continued engineering high-throughput runtime 
C++ systems, building an Unreal Engine C++ plugin that captures and exports runtime character trajectory, 
pose, and velocity state to drive motion-matching locomotion query databases. Beyond my studio work 
at EA, I have engineered a custom C++ game engine (\emph{Hoowan Game Engine}) featuring an 
entity-component system, input polling, and collision resolution pipeline, as well as a 3D 
first-person shooter (\emph{DoomSlop}) with responsive movement and camera controls, 
hitscan weapon raycasting, combat state feedback, and a custom radar minimap.

As an internal EA engineer based in the Pacific Northwest with a deep appreciation for Apex Legends' 
industry-leading movement and combat feel, I am excited about the opportunity to collaborate with 
Respawn's team to deliver polished, reliable features for a live multiplayer audience. 
Thank you for your time and consideration, and I look forward to discussing how my experience can 
contribute to the team.

\vspace{10pt}
\noindent Sincerely,

\vspace{14pt}
\noindent \textbf{Juan Becerra}

\end{document}
